\section{Die Bibel das Wort Gottes}

Eigentlich wollte ich eine Serie zum Thema Gemeinde beginnen. Wir haben ja bis jetzt die Heilsgeschichte, der Rote Faden der Bibel, angeschaut und danach, die Namen Gottes in der Bibel. Egal welches theologische Thema ich anfange, ohne die Bibel, würde ich hier vorne stehen und euch irgendwelche erfundenen Geschichten erzählen. Und glaubt mir, ich bin alles andere als ein guter Geschichtenerzähler. Darum habe ich alles aus der Bibel und aus Büchern, die diese Themen aus der Bibel aufgreifen und analysieren.

Darum habe ich das Thema Gemeinde auf die Seite gelegt und fange nun das Thema über die Bibel an. Anfangs September habe ich einen Vortrag, bei dem ich die Entstehung und die Einzigartigkeit der Bibel aufzeige. Bei diesem Vortrag ist es mir wichtig zu zeigen, das die Bibel nicht einfach ein Buch ist, wie jedes andere und er die Leute anregen soll, darüber nachzudenken.

Hier bei dieser Serie, gehe ich ein Schritt weiter. Ich denke, alle die hier sitzen, glauben wie auch ich, dass die Bibel Gottes Wort ist. Aber was heisst das, Gottes Wort? Ist sie vom Himmel gefallen? Hat er die Bibel den Schreiben diktiert? Haben die Schreiber geträumt und das geträumte aufgeschrieben, wie Paulus in 2. Korinther sagt, dass er nicht weiss ob wer wirklich im Himmel oder nur geträumt hat? \biblezit{IIKor}{12:2}? Was heisst von Gott inspiriert? 

Ein weiters Thema ist auch, wie stimmen die heutigen Übersetzungen? Wie können wir sicher sein, dass das was hier in dieser Schlachter steht auch wirklich das Wort Gottes ist? Warum gibt es so viele Übersetzungen? Welche Übersetzungen sind denn die besten? Was ist der Unterschied zwischen Schlachter und Elberfelder?

Aus diesen Fragen ergibt sich ein Ablauf mit folgenden Inhalt:
\begin{enumerate}
    \item Einführung in die Schrift
    \item Die inspiration der Schrift
    \item Die Autorität der Schrift
    \item Die Irrtumslosigkeit der Schrift
    \item Die Bewahrung der Schrift
    \item Das Lehren und Predigen der Schrift
    \item Die Verpflichtung gegenüber der SChrift
\end{enumerate}

Dieses Inhaltsverzeichnis, werde ich nächstes mal auf dem Beamer zeigen, kommt aus dem grossen Theologie Buch von John MacArthur. Ich werde mich ein bisschen an dem Ablauf in diesem Buch orientieren und vieles wird auch aus den Vorträgen von Roger Liebi kommen. 

John McArthur schreibt:
\begin{quote}
    Die Lehre von der Schrift ist absolut fundamental und essenziell, da sie die einzige wahrhaftige Quelle aller christlichen Wahrheit ist.
\end{quote}

Dieser Satz ist der Grund, wieso ich dieses Thema gegenüber der Gemeinde vorgezogen habe. Die Bibel ist nun mal das einzige was wir haben, aus der wir unserer Glaube und uns kontrollieren können.  Wir sind  eine sogenanne Schriftreligion. 2000 Jahre nachdem Jesus gestorben ist, ist dieses Buch das einzige was wir haben. Darum sagt auch Jesus in \biblezit{Joh}{20:29}. Jesus weiss wie schwer wir uns mit dem Glauben tun. Das dieses Buch noch existiert, ist kein Zufall. Das erkläre ich dann am 3. September.

Es reicht wenn wir uns an die Bibel halten. Die Mormonen und Zeugen Jehowas haben auch die Bibel. 

Aber wir müssen immer aufpassen. Die Mormonen haben die Bibel plus das Buch Mormon. Die Zeugen Jehowas haben die Bibel plus den Wachturm. Die orthodoxen Juden haben die Bibel plus die Auslegungsbücher \zB{} der Talmut. Die katholische Kirche hat die Bibel plus die Aussagen der Kirchenväter und der Katechismus.\\
Es ist immer die Bibel plus. Auch wir müssen aufpassen. Vor allem heute im Zeitalter der sozialen Medien. Es gibt soviele Experten, dass man nicht mehr weiss, was stimmt und was nicht.

Es gibt super Predigten vom Mitternachtsruf, von Roger Liebi, John McArthur und vielen anderen, aber die sind nicht der Massstab, besser gesagt nicht unsere Grundlage des Glauben. Sie können uns helfen die Bibel besser zu verstehen. Aber es muss nicht das richtige sein. Das einzig richtige steht in der Bibel. Darum ist es so wichtig, dass wir diese selber lesen. Nur so können wir prüfen, dass wir auf dem Weg Jesus Christus sind.

Alle vorher erwähnten Sekten oder Glaubensgruppen, wollen nicht das ihre Schäfchen die Bibel selber lesen. Zeugen Jehowas haben ihr Leitungsgremium das die Bibel auslegt. Die Mormonen die Heiligen der letzten Tage. Die katholische Kirche ihre Pfarrer und Bischöfe.\\
Ich habe einem Katholiken mal geschrieben, dass die Bibel eigentlich einfach zu verstehen ist und wir d
ie gut selber lesen können. Er schrieb zurück, dass sogar Petrus geschrieben hat, dass die Bibel schwer zu verstehen ist und es darum Leute braucht die studiert haben um uns die Bibel zu erklären. Ich fragte ihn dann, ob er denn die Stelle zu Ende gelesen hätte, weil da steht, das genau diese Ausleger das Problem sind.

Ich will heute nicht länger werden und dann in zwei Wochen mit dem Thema Bibel anfangen.